\section*{ОБОЗНАЧЕНИЯ И СОКРАЩЕНИЯ}

БД -- база данных.

ИС -- информационная система.

ИТ -- информационные технологии. 

КТС -- комплекс технических средств.

ПО -- программное обеспечение.

СУБД -- система управления базами данных.

ТЗ -- техническое задание.

ТП -- технический проект.

API (Application Programming Interface) -- программный интерфейс приложения.

CI/CD (Continuous Integration/Continuous Delivery) -- непрерывная интеграция и доставка изменений программного обеспечения.

Docker -- платформа контейнеризации приложений.

Grafana -- система визуализации метрик и построения мониторинговых панелей.

HTTP (Hypertext Transfer Protocol) -- протокол передачи данных в веб-приложениях.

JSON (JavaScript Object Notation) -- текстовый формат обмена структурированными данными.

JWT (JSON Web Token) -- формат токена для передачи утверждений между участниками системы.

Kafka -- распределенная платформа обмена событиями и потоковой обработки данных.

Ktor -- фреймворк для разработки серверных приложений на языке Kotlin.

Kubernetes -- платформа оркестрации контейнеризированных приложений.

OAuth -- протокол авторизации для предоставления ограниченного доступа к ресурсам.

Prometheus -- система мониторинга и сбора метрик.

React -- библиотека для разработки пользовательских интерфейсов.

REST (Representational State Transfer) -- архитектурный стиль построения веб-API.

TypeScript -- язык программирования на основе JavaScript со статической типизацией.

UML (Unified Modelling Language) -- язык графического описания для объектного моделирования в области разработки программного обеспечения.
